\documentclass[pdflatex,sn-mathphys-num]{sn-jnl}

\usepackage{graphicx}
\usepackage{url}
\usepackage{amsmath,amssymb,amsfonts}
\usepackage{amsthm}
\usepackage[title]{appendix}
\usepackage{booktabs}

\theoremstyle{thmstyleone}
\newtheorem{theorem}{Theorem}
\newtheorem{proposition}[theorem]{Proposition}
\newtheorem{lemma}[theorem]{Lemma}
\newtheorem{corollary}[theorem]{Corollary}
\newtheorem{conjecture}[theorem]{Conjecture}

\theoremstyle{thmstyletwo}
\newtheorem{example}{Example}
\newtheorem{remark}{Remark}

\theoremstyle{thmstylethree}
\newtheorem{definition}{Definition}

\newcommand{\evpi}{\mathrm{EVPI}}
\newcommand{\evppi}{\mathrm{EVPPI}}
\newcommand{\cvar}{\mathrm{CVaR}}
\newcommand{\ubound}{U}

\raggedbottom

\begin{document}

\title[Decision Value of Hidden Supply-Network Structure]{The Decision Value
of Hidden Supply-Network Structure: Zero-Value Conditions and
Identification-Premium Bounds}

\author*[1]{\fnm{Rahul Kumar} \sur{Mishra}}\email{mishra.rahul98@gmail.com}

\affil*[1]{\orgname{Independent Researcher}, \orgaddress{\city{Ranchi},
\state{Jharkhand}, \country{India}}}

\abstract{Firms invest in sub-tier supply-network visibility through
mapping, supplier questionnaires and audits, and a substantial literature
studies how hidden supplier structure can be inferred. We study a prior
question: what is that structure worth to a specific mitigation decision? We
formulate discovery-coupled mitigation as a probe-then-secure problem, with
no probe cost priced here, on a partially observed supply directed acyclic
graph and analyse the expected value of perfect structural information
(EVPI). We give sufficient structural conditions, namely invariance of the
loss distribution under identification-free plans together with uniform
plan-level dominance of identification-free actions, under which EVPI is
exactly zero; since perfect information dominates any acquisition policy,
every feasible information-acquisition policy then has zero decision value
within the modelled problem. The conditions are stated for any monotone
law-invariant objective, so they cover expected loss and conditional value
at risk alike. The first condition holds when a symmetry of the public
network exchanges the candidate suppliers of each category together with
their consumers, upstream subtrees and correlation clusters; exchangeability
of the candidates' marginal risks alone does not suffice, and we give a
counterexample. For the second we give a sufficient condition read off the
graph and the cost vector. Outside the zero-value regime we bound EVPI above
by a sum of per-category identification premiums computable from prior and
public data, prove the bound for node-disjoint categories, give the premium
exactly in a dilution special case, and identify a sub-case in which the
bound is attained. Exact enumeration of synthetic miniature instances found
no violation of any stated result. The conditions and the bound are not
evaluated on real supply networks here.}

\keywords{value of information, EVPI, supply network visibility, disruption
risk, conditional value at risk, endogenous uncertainty}

\maketitle

\section{Introduction}\label{sec:intro}

Sub-tier supply-network visibility is bought through mapping services,
supplier questionnaires, audit programmes and related discovery activities.
The case for such spending is usually argued from the operational
consequences of incomplete visibility. It is much
less common to price the information itself: to ask by how much a firm's
mitigation decisions would improve if the hidden structure were supplied at
no cost.

Decision analysis answers that question with the expected value of perfect
information, which bounds the improvement obtainable from any amount of
investigation. Applying it to supply-network structure is easy to state and
harder to use, because the decision the information feeds --- which
suppliers to qualify a backup for, under a cost budget, on a network whose
upstream shape is the unknown --- is combinatorial.

We formulate this problem as a decision model of discovery-coupled
mitigation and characterise its information value. One point of scope should
be clear from the outset: we value the information, not the coupling. No
probe cost and no shared probe-and-mitigation budget enters any quantity
studied here, since the uninformed planner acts on the prior and the
informed planner is handed $\theta$ outright. The first result is
negative and, in our view, the more useful one: there are structural
conditions under which the value of perfect structural information is exactly
zero. Both can be examined before any structural discovery is commissioned: the
first, a symmetry of the public network that exchanges the candidates of each
category together with their consumers, upstream subtrees and correlation
clusters, directly from public data, and the second through a sufficient condition on the graph
and the cost vector. The result is independent of
the optimisation algorithm used to implement the decision: under those
conditions no information-acquisition policy can improve the objective,
because perfect information dominates every such policy and is itself worth
nothing. The second result is a bound for the cases where those conditions
fail, expressed as a sum of per-category identification premiums, each
computable from information the firm already holds.

\paragraph{Central claim.} We characterise sufficient structural conditions
under which the decision value of hidden supplier information is exactly
zero, and derive a computable upper bound on that value outside the
zero-value regime. Everything else in the paper supports that claim: the
enumeration experiments are verification, the CVaR statement is an
extension, and the shared-subtree case is an unresolved extension we label as
such.

\paragraph{Contributions.}
\begin{enumerate}
\item \emph{Zero-value characterisation.} Theorem~\ref{thm:zero} gives
sufficient conditions for $\evpi=0$, and Lemma~\ref{lem:garbling} transfers
the conclusion to every information-acquisition policy.
Lemma~\ref{lem:H1suff} gives a sufficient condition for the first
hypothesis, a symmetry of the public structure, and
Example~\ref{ex:marginal} shows that exchangeability of the candidates'
marginal risks does not suffice. Proposition~\ref{prop:witness} gives a
sufficient condition for the harder of the two hypotheses that is read off
the public graph and cost vector.
Corollary~\ref{cor:lawinv} extends the statement to monotone law-invariant
objectives, covering conditional value at risk.
\item \emph{Identification-premium upper bound.} Theorem~\ref{thm:premium}
bounds EVPI by a finite, directly computable sum of per-category premiums
for node-disjoint categories, and Corollary~\ref{cor:evppi} bounds the
partial value of each category separately. Computable EVPI bounds are not new
(Section~\ref{sec:related}, stream A$'$). What we believe to be new, and have
not found stated elsewhere, is a bound of this form: it requires no convexity,
decomposes over a discrete latent assignment, and is localised to individual
channels: under the assumptions of Theorem~\ref{thm:premium}, a channel
whose premium is zero has zero partial value by Corollary~\ref{cor:evppi}.
Whether ranking channels by premium identifies those worth discovering is
not established here.
Further, Corollary~\ref{cor:dilution} gives the exact
premium in a dilution special case, and we identify a sub-case of it in
which the bound is attained with equality; and
Proposition~\ref{prop:disp} gives the premium in closed form under a
specified dispersion family, identifying a threshold below which dispersion
does not enter the bound at all.
\item \emph{Computational verification.} Exhaustive enumeration of exact
miniature instances tests the implementation of the definitions, compares
the per-channel premium with exact EVPPI, and searches for violations of the
stated results and for counterexamples to Conjecture~\ref{conj:shared}
(Section~\ref{sec:verify}). It is evidence about small synthetic instances,
not a measurement on supply networks.
\end{enumerate}

\noindent Section~\ref{sec:limits} states what the paper does not establish.

\section{Related work and positioning}\label{sec:related}

\paragraph{(A) Value of information.} The expected value of perfect
information originates with Raiffa and Schlaifer~\cite{raiffa1961} and
Howard~\cite{howard1966}. Its partial form, EVPPI, is standard in health
economics, where the computational cost of nested estimation is a recognised
difficulty~\cite{strong2014,heath2024}. Our premium is an EVPPI-shaped
quantity for structural parameters of a network, evaluated by finite
enumeration over budget breakpoints rather than by nested simulation. We
claim no priority over that literature, whose machinery we reuse.

\paragraph{(A$'$) Bounding EVPI when it cannot be computed.} We must be
explicit that a computable upper bound on EVPI is not a new idea. Because
exact evaluation is often prohibitive, bounds have been derived since
Madansky~\cite{madansky1960}, and Avriel and Williams~\cite{avriel1970}
obtained EVPI bounds evaluable as part of a post-optimal analysis of a linear
or convex program. Ziemba and Butterworth~\cite{ziemba1975} and Hausch and
Ziemba~\cite{hausch1983} devote a two-part study to bounding the value of
information directly; Huang, Vertinsky and Ziemba~\cite{huang1977} sharpen
the perfect-information bounds for simple-recourse problems; and
Birge~\cite{birge1982} develops the companion bounds for the value of the
stochastic solution. That family of
bounds rests on convexity or concavity: Huang, Vertinsky and Ziemba state
theirs for simple-recourse problems with nonlinear return functions and
concave utility functions, using extensions of the Jensen and
Edmundson--Madansky inequalities. Theorem~\ref{thm:premium} uses no
convexity and no continuity: it is a combinatorial decomposition over a
discrete latent assignment, evaluated on a finite set of budget breakpoints,
and it localises the bound to individual categories, so that a category whose
premium is zero is shown to have zero partial value. Our contribution is that
specific form, not the general idea of bounding EVPI, and where the two apply to
the same problem we make no claim that ours is tighter.

\paragraph{(B) Information acquisition and active learning.} Adaptive
information gathering with performance guarantees is well
developed~\cite{golovin2011}, and decision-coupled acquisition, which drives
uncertainty into a single decision region instead of reducing it uniformly,
is the closest precedent for the objective used here~\cite{javdani2014}. The ordering of information structures by their
decision value is Blackwell's~\cite{blackwell1953}, and
Lemma~\ref{lem:garbling} is a standard consequence of it rather than a
contribution of this paper. That literature asks how to acquire information
efficiently. We ask when acquisition cannot help at all, which is logically
prior and is not settled by any performance guarantee.

\paragraph{(C) Decision-dependent information discovery.} Stochastic
programming with endogenous uncertainty studies problems in which decisions
determine what is revealed~\cite{jonsbraten1998,goel2006}, and the robust
counterpart optimises under decision-dependent
revelation~\cite{vayanos2025}. Artstein and Wets~\cite{artstein1993} and
Artstein~\cite{artstein1994,artstein1999} treat the information structure
itself as a decision variable through the notion of sensors;
Artstein~\cite{artstein1994} recasts a two-stage problem with a probing
decision as a three-stage stochastic program and establishes a strong law of
large numbers for it, and Artstein~\cite{artstein1999} addresses the
trade-off between the gain from new information and the cost of obtaining
it. Closest to our setting is the probing-enhanced
stochastic program of Ma et al.~\cite{ma2024}, in which costly discrete
probes are chosen before a two-stage decision and the perfect-information
bound is the benchmark to be improved upon; their contribution is
algorithmic, a branch-and-bound scheme over probing sets. The direction of the question differs: that stream computes policies
\emph{under} revelation, while we characterise conditions in which revelation
changes nothing, so that the
problem those methods solve has a degenerate answer.

\paragraph{(D) Valuation of supply-chain visibility.} It would be wrong to
claim that visibility has not been valued. Caridi et al.~\cite{caridi2014}
develop a structured value-assessment model for visibility and apply it to a
real chain, and recent student and practitioner work continues to estimate
returns to deeper tiers~\cite{cermeno2026}. That work maps visibility improvements to
operational and strategic benefits, typically through causal models of
performance indicators. Our question is narrower and its answer sharper: for
a specified mitigation decision and a specified loss, when is the decision
value of latent structure exactly zero, and when it is not, what upper bound
can be computed in advance. Neither statement subsumes the other.

\paragraph{(E) Multi-tier structure and hidden links.} Inference of hidden
supplier relationships is an active area~\cite{brintrup2018,kosasih2022},
and topological measures of network vulnerability are well
established~\cite{wagner2010,simchi2015}. These methods find or score
structure; they do not price it against a decision.

\paragraph{(F) Disruption mitigation.} Mitigation and contingency strategy
follows Tomlin~\cite{tomlin2006}, with the wider OR/MS literature surveyed
by Snyder et al.~\cite{snyder2016}. Our mitigation model --- budgeted backup
qualification under a conjunctive cascade --- is deliberately simple, since
the object of study is the information rather than the mitigation policy.

\paragraph{(G) Risk measures.} The tail machinery is that of Rockafellar and
Uryasev~\cite{rockafellar2000} and Artzner et al.~\cite{artzner1999}.

Table~\ref{tab:novelty} summarises where this paper sits against these
streams. The column that separates it is the third.

\begin{table}[ht]
\caption{Positioning against adjacent streams. The column that separates
this paper is the sufficient zero-value condition; computable EVPI bounds
already exist (stream A$'$), and our claim there concerns the form of the
bound rather than its existence. The zero-value column records what we found
in the works we examined, not an exhaustive claim about each
stream. ``Model only'' means that our results are stated for a model of a
supply network and are not evaluated on real networks.}\label{tab:novelty}
\small
\begin{tabular}{@{}p{3.7cm}cccp{1.8cm}c@{}}
\toprule
& values & gives a & gives an & bound is & applied to \\
stream & latent & zero-value & EVPI & convexity- & supply \\
& structure & condition & bound & free & networks \\
\midrule
Value of information (A) & yes & no & no & n/a & no \\
EVPI bounding (A$'$) & no & no & yes & not in those examined & no \\
Active acquisition (B) & yes & no & no & n/a & no \\
Decision-dependent discovery (C) & yes & no & no & n/a & no \\
Visibility valuation (D) & yes & no & no & n/a & yes \\
Hidden-link inference (E) & no & no & no & n/a & yes \\
This paper & yes & yes & yes & yes & model only \\
\bottomrule
\end{tabular}
\end{table}

\section{Model}\label{sec:model}

\paragraph{Network.} A finite directed acyclic graph $G=(V,E)$, where
$u\to v$ means $u$ supplies $v$. A focal node $0$ consumes $K_0$ input
categories. Every node $v$ consumes categories $k\in K_v$, each served by a
set of suppliers. Service is conjunctive across categories and substitutable
within a category, in the sense that any one working supplier suffices. That
clause fixes the semantics of the cascade rather than the sourcing structure.
The candidate set is epistemic rather than a redundancy: exactly one candidate
actually supplies the category, and the set records who it might be. The
model therefore contains no concurrently multi-sourced category, and dual
sourcing as a mitigation lies outside it. Each node has a publicly known securing cost $c_v$ and
belongs to a publicly known correlation cluster. A node's \emph{tier} is its
distance from the focal node along the supply direction: the focal firm's
direct suppliers are tier one, their suppliers tier two, and so on.

\paragraph{Uncertainty.} The structure parameter $\theta$ specifies, for
each latent category $(v,k)$, which candidate in a finite set $C_{v,k}$ is
the true supplier, exactly one candidate serving the category, an assumption
Corollary~\ref{cor:dilution} and Proposition~\ref{prop:disp} use and, in the
dispersion family of Definition~\ref{def:disp}, which candidate carries which
risk attribute. A
prior $\pi$ over $\theta$ is given. Probing a node reveals its candidate
in-edges through a noisy asymmetric channel. That channel belongs to the
wider model, of which this paper studies only the value of the information a
probe could deliver; it enters no quantity studied here, since the
uninformed planner sees only the prior and the informed planner sees
$\theta$ outright.

\paragraph{Disruption and loss.} A scenario $\omega$ draws node-level
failures, correlated within clusters by a one-factor Gaussian copula.
Operability propagates through the conjunctive cascade. A plan
$S\subseteq V\setminus\{0\}$ with $c(S)\le B$ qualifies backup sources:
$u\in S$ guarantees $u$'s output to its consumers. The focal loss is the
fraction of focal categories starved,
\begin{equation}
L(S;\theta,\omega)=\frac{1}{K_0}\sum_{k=1}^{K_0} Z_k,
\qquad Z_k=\mathbf{1}\{\text{focal category }k\text{ starved}\}.
\label{eq:loss}
\end{equation}

\paragraph{Objectives and the exact target.} Write
$j(S;\theta)=\mathbb{E}_\omega L(S;\theta,\omega)$ and
$J(S;\theta)=\cvar_{\alpha,\omega}(L(S;\theta,\omega))$. Let $S_\pi$
minimise $\mathbb{E}_{\theta\sim\pi}j(S;\theta)$ and $S_\theta$ minimise
$j(S;\theta)$. Then
\begin{equation}
\evpi(\pi)=\mathbb{E}_\pi\!\left[j(S_\pi;\theta)\right]
-\mathbb{E}_\pi\!\left[j(S_\theta;\theta)\right]\;\ge\;0 .
\label{eq:evpi}
\end{equation}

\paragraph{Units.} $L$ is a fraction of focal categories, so every quantity
here is dimensionless. A value of $0.005$ for $\evpi$ or for the bound below
means that perfect structural knowledge would reduce the objective by half a
percentage point of the focal basket: expected category starvation where the
objective is $j$, and the conditional tail average of the starved fraction
where it is $J$. Because $L$ is normalised by $K_0$, rescuing one focal
category moves the objective by $1/K_0$, so the same absolute value means
different things on networks with different numbers of focal categories, and
$K_0$ times a value counts the focal categories rescued. Converting either
into money requires the cost of one starved category, which the model does
not represent.

\paragraph{Channels.} A \emph{channel} is one latent category at one node, and
is the unit at which identification premiums are computed. A \emph{focal
category} is one of the $K_0$ inputs the focal node consumes.

\subsection{What an implemented planner observes}

A firm operates an implemented planner rather than an exact optimiser, and
what it can observe by running that planner with and without $\theta$ is
not \eqref{eq:evpi}. Let
$\hat S_\pi$ and $\hat S_\theta$ be the plans returned by the implemented
planner given the prior and given $\theta$. Define the \emph{heuristic
oracle advantage}
\begin{equation}
H \;=\; \mathbb{E}_\pi\!\left[j(\hat S_\pi;\theta)\right]
      -\mathbb{E}_\pi\!\left[j(\hat S_\theta;\theta)\right].
\label{eq:H}
\end{equation}

\begin{proposition}[Exact relation between $H$ and EVPI]\label{prop:H}
Let $\varepsilon_\pi=\mathbb{E}_\pi[j(\hat S_\pi;\theta)]
-\mathbb{E}_\pi[j(S_\pi;\theta)]\ge0$ and
$\varepsilon_\theta=\mathbb{E}_\pi[j(\hat S_\theta;\theta)
-j(S_\theta;\theta)]\ge0$ be the planner's optimality gaps in the two
regimes. Then
\begin{equation}
H=\evpi(\pi)+\varepsilon_\pi-\varepsilon_\theta .
\label{eq:Hrelation}
\end{equation}
Hence $H\le\evpi$ if and only if $\varepsilon_\pi\le\varepsilon_\theta$. In
general the sign of $\varepsilon_\pi-\varepsilon_\theta$ is not determined,
so $H$ is neither an upper nor a lower bound on EVPI; if the planner is
$\theta$-optimal when given $\theta$, so that $\varepsilon_\theta=0$, then
$H\ge\evpi$.
\end{proposition}

\begin{proof}
(E1) By definition of $\varepsilon_\pi$,
$\mathbb{E}[j(\hat S_\pi;\theta)]=\mathbb{E}[j(S_\pi;\theta)]
+\varepsilon_\pi$.

\noindent(E2) By definition of $\varepsilon_\theta$,
$\mathbb{E}[j(\hat S_\theta;\theta)]=\mathbb{E}[j(S_\theta;\theta)]
+\varepsilon_\theta$.

\noindent(E3) Subtracting (E2) from (E1) and using \eqref{eq:evpi} gives
\eqref{eq:Hrelation}. Both gaps are nonnegative because $S_\pi$ and
$S_\theta$ minimise over the same feasible sets that the planner searches
heuristically, and neither gap dominates the other in general, which gives
the remaining statements.
\end{proof}

Proposition~\ref{prop:H} holds for any fixed objective used in every term;
under the CVaR objective $J$ it relates $H$ to the CVaR value of perfect
information. A planner's observed advantage is therefore a planner-specific
quantity, and it bounds $\evpi$ from neither side unless the sign of
$\varepsilon_\pi-\varepsilon_\theta$ is known.

\section{When structural information has zero value}\label{sec:zero}

\begin{lemma}[Acquisition is dominated by perfect information]
\label{lem:garbling}
For any information-acquisition policy $\varphi$ with value
$\mathrm{Val}(\varphi)$, $\mathrm{Val}(\varphi)\le\evpi(\pi)$.
\end{lemma}

\begin{proof}
(L1) By definition,
$\mathrm{Val}(\varphi)=\mathbb{E}_\pi[j(S_\pi;\theta)]
-\mathbb{E}_o\mathbb{E}_{\theta\mid o}[j(S_{\pi_o};\theta)]$, where $o$ is
the observation induced by $\varphi$, $\pi_o$ the posterior and $S_{\pi_o}$
the plan chosen after $o$.

\noindent(L2) Pointwise in $(o,\theta)$, $S_{\pi_o}$ is one feasible plan
among others, so $j(S_{\pi_o};\theta)\ge\min_S j(S;\theta)
=j(S_\theta;\theta)$.

\noindent(L3) Taking $\mathbb{E}_{\theta\mid o}$ of (L2), then
$\mathbb{E}_o$, and applying the tower property gives
$\mathbb{E}_o\mathbb{E}_{\theta\mid o}[j(S_{\pi_o};\theta)]
\ge\mathbb{E}_\theta[j(S_\theta;\theta)]$.

\noindent(L4) Substituting (L3) into (L1) yields
$\mathrm{Val}(\varphi)\le\evpi(\pi)$.
\end{proof}

\begin{proposition}[Bernoulli CVaR]\label{prop:bern}
For $Z\in\{0,1\}$ with $\mathbb{P}(Z=1)=p$,
$\cvar_\alpha(Z)=\min\{1,p/(1-\alpha)\}$, and
$\cvar_\alpha(K_0^{-1}\sum_kZ_k)\le
K_0^{-1}\sum_k\min\{1,p_k/(1-\alpha)\}$.
\end{proposition}

\begin{proof}
(P1) If $p\ge1-\alpha$ the worst $(1-\alpha)$ tail consists of ones only, so
the tail average is $1$.

\noindent(P2) If $p<1-\alpha$ the tail contains all $p$ mass of ones and
$(1-\alpha-p)$ of zeros, so its average is $p/(1-\alpha)$.

\noindent(P3) Combining gives $\cvar_\alpha(Z)=\min\{1,p/(1-\alpha)\}$.

\noindent(P4) CVaR is subadditive and positively
homogeneous~\cite{artzner1999}, so
$\cvar_\alpha(K_0^{-1}\sum_kZ_k)\le K_0^{-1}\sum_k\cvar_\alpha(Z_k)$, and
(P3) gives the stated form.
\end{proof}

\begin{definition}[Identification-free actions]\label{def:free}
An action is \emph{identification-free} for the latent category $(v,k)$ if
the conditional law of the focal loss $L$, given a plan containing that
action, does not depend on $\theta_{v,k}$. A plan is
\emph{identification-free} if each of its elements is identification-free for
the category it addresses. This is a condition on the plan's own elements
only: it says nothing about the latent categories the plan leaves untouched,
whose $\theta$-dependence is what (H1) below rules out.
Canonical examples of identification-free actions:
secure the consumer $v$;
secure every candidate in $C_{v,k}$; and, under the sufficient condition
(H1$'$) below, take no action, since every structure automorphism fixes the
empty plan. Under heterogeneous latent candidate risk the last of these is not
identification-free, and exchangeability of the candidates' marginal laws
alone does not make it so, since the law of $L$ under inaction can depend on which
candidate is the supplier. Securing one named
candidate $u\in C_{v,k}$ is \emph{identification-dependent}, since its
effect depends on whether $u$ is the supplier.
\end{definition}

We use two hypotheses.
\begin{itemize}
\item[(H1)] \emph{Loss-law invariance.} For every identification-free plan
$S$, the law of $L(S;\theta,\cdot)$ is the same for all $\theta$. The
hypothesis is not implied by Definition~\ref{def:free}: an
identification-free plan constrains only the categories its own elements
address, and (H1) is exactly the requirement that the categories it leaves
alone contribute a law that does not vary with $\theta$. That is why the
sufficient condition (H1$'$) below is stated over every latent category rather
than over the plan's.
\item[(H2)] \emph{Uniform plan-level dominance.} Every
identification-dependent element $u$ of any plan admits an
identification-free replacement $a(u)$ with $c(a(u))\le c(u)$ and
$j(S\setminus\{u\}\cup\{a(u)\};\theta)\le j(S;\theta)$ for all $\theta$.
Securing every candidate in the category is identification-free but cannot
serve as $a(u)$ under positive costs: the replacement set contains $u$, so it
is never cheaper. Only an action at the consumer level, or outside the
category, can witness the cost clause.
\end{itemize}

\paragraph{A checkable sufficient condition for (H1).} Exchangeability of the
candidates of each category in marginal operability law, in securing cost and
under the prior is \emph{not} sufficient. Two candidates with the same failure
law can differ in their other consumers, in their upstream subtrees or in their
correlation cluster, and each of these differences can make the joint law of
the starvation indicators, and with it the law of $L$, depend on which
candidate is the supplier. Example~\ref{ex:marginal} in
Appendix~\ref{app:H1} gives an instance in which the marginal condition holds,
(H2) holds by Proposition~\ref{prop:witness}, and $\evpi>0$. What suffices is
a symmetry of the public structure. Write $q_v$ for the marginal failure
probability of node $v$ and $\kappa(v)$ for its cluster label. A
\emph{structure automorphism} is a bijection $\sigma:V\to V$, together with
bijections $\sigma_v:K_v\to K_{\sigma v}$ and a bijection $\eta$ of cluster
labels, such that
(A0) $\sigma(0)=0$;
(A1) whenever $(v,k)$ has public supplier $s$, $(\sigma v,\sigma_vk)$ has
public supplier $\sigma(s)$;
(A2) $(v,k)$ is latent if and only if $(\sigma v,\sigma_vk)$ is, and then
$C_{\sigma v,\sigma_vk}=\sigma(C_{v,k})$;
(A3) $q_{\sigma v}=q_v$, $\kappa(\sigma v)=\eta(\kappa(v))$, and $\eta$
preserves each cluster's copula parameter.
Because $\sigma$ preserves every supply and candidate relation, it carries
each node's upstream subtree and consumers along with the node. It acts on
states by $(\sigma\theta)_{\sigma v,\sigma_vk}=\sigma(\theta_{v,k})$ and on
plans by $\sigma S=\{\sigma u:u\in S\}$.
\begin{itemize}
\item[(H1$'$)] \emph{Automorphic exchangeability.} For every
identification-free plan $S$ and all $\theta,\theta'$ there is a structure
automorphism $\sigma$ with $\sigma\theta=\theta'$ and $\sigma S=S$.
\end{itemize}
Lemma~\ref{lem:H1suff} in Appendix~\ref{app:H1} proves that (H1$'$) implies
(H1); in fact the joint law of $(Z_1,\dots,Z_{K_0})$ is then invariant in
$\theta$ up to a fixed permutation of the focal categories. Neither the prior
nor the cost vector enters, since (H1) concerns a fixed plan pointwise in
$\theta$. A local form is easier to check. It suffices that for every latent
category $(v,k)$ and all $u,u'\in C_{v,k}$ some structure automorphism maps
$u$ to $u'$, fixes $v$ and $k$, fixes pointwise every other candidate set and
every other latent category's consumer, and fixes every identification-free
plan. The typical witness swaps the upstream cones of two candidates when the
cones are disjoint, isomorphic with matching failure probabilities and cluster
labels, and joined to the rest of the network only by the edge into $(v,k)$.
A candidate that belongs to two candidate sets, or a candidate with a
consumer that its rivals lack, defeats the local form.

\begin{theorem}[Zero decision value of hidden structure]\label{thm:zero}
Under (H1) and (H2), $\evpi(\pi)=0$; hence by Lemma~\ref{lem:garbling} every
feasible information-acquisition policy has zero decision value.
\end{theorem}

\begin{proof}
(T1) Fix $\theta$. The feasible plan set is finite, so a $\theta$-optimal
plan $S_\theta$ exists.

\noindent(T2) Suppose $S_\theta$ contains an identification-dependent
element $u$, and replace it by $a(u)$. Feasibility is preserved because
$c(a(u))\le c(u)$, and the objective does not increase; both by (H2).

\noindent(T3) Each such replacement strictly decreases the number of
identification-dependent elements, because $a(u)$ is identification-free by
(H2) and the status of no other element changes. The plan is finite, so the
iteration terminates at an identification-free plan $\tilde S_\theta$ that
is feasible and satisfies
$j(\tilde S_\theta;\theta)\le j(S_\theta;\theta)$.

\noindent(T4) By $\theta$-optimality of $S_\theta$ the reverse inequality
also holds, so $j(\tilde S_\theta;\theta)=j(S_\theta;\theta)$.

\noindent(T5) By (H1) the value of any identification-free plan does not
depend on $\theta$. Let $j^\star$ be the minimum over identification-free
plans, attained at $S^\star$; both are $\theta$-independent.

\noindent(T6) Since $\tilde S_\theta$ is identification-free,
$j(\tilde S_\theta;\theta)\ge j^\star$. Since $S^\star$ is feasible,
$j^\star=j(S^\star;\theta)\ge j(S_\theta;\theta)=j(\tilde S_\theta;\theta)$.
The two inequalities force $j(S_\theta;\theta)=j^\star$ for every $\theta$.

\noindent(T7) Hence $\mathbb{E}_\theta[j(S_\theta;\theta)]=j^\star$.

\noindent(T8) $S^\star$ is available to the prior planner, whose optimum can
only be at least as good:
$\mathbb{E}_\theta[j(S_\pi;\theta)]\le j^\star$.

\noindent(T9) Substituting (T7) and (T8) into \eqref{eq:evpi} gives
$\evpi\le0$; with $\evpi\ge0$ this forces $\evpi=0$.
\end{proof}

The argument requires both hypotheses. With (H2) alone, identification-free
plans are optimal but their value may still depend on $\theta$, and step (T5)
fails. The argument uses nothing specific to expected loss, which gives the
following extension.

\begin{corollary}[Monotone law-invariant objectives]\label{cor:lawinv}
Let $\rho$ be a functional of the loss distribution that is monotone with
respect to almost sure dominance, and set
$J_\rho(S;\theta)=\rho(L(S;\theta,\cdot))$. If (H1) holds and (H2) holds
with $j$ replaced by $J_\rho$, then the value of perfect structural
information under $J_\rho$ is zero. This covers expected loss,
$\cvar_\alpha$ for every $\alpha$, and spectral risk measures.
\end{corollary}

\begin{proof}
(C1) Steps (T1)--(T4) use only finiteness of the plan set, the cost clause
and the dominance inequality, all assumed for $J_\rho$.

\noindent(C2) Step (T5) uses (H1) through the \emph{law} of $L$: if the law
of $L(S;\theta,\cdot)$ does not depend on $\theta$, then neither does
$\rho(L(S;\theta,\cdot))$ for any functional $\rho$ of that law. This is
where law invariance rather than mean invariance is required.

\noindent(C3) Steps (T6)--(T9) are algebraic and carry over with $j$
replaced by $J_\rho$.
\end{proof}

\section{Bounding the value outside the zero-value regime}\label{sec:bound}

\begin{theorem}[Identification-premium bound, disjoint categories]
\label{thm:premium}
Suppose the focal categories' subtrees are node-disjoint, the $\theta_k$ are
independent across categories, and plans decompose with
$c(S)=\sum_k c(S_k)$. For a budget $b$ define
\begin{align}
v_k^{\mathrm{inv}}(b)&=\min_{c(a)\le b}\;
\mathbb{E}_{\theta_k}\,p_k(a;\theta_k),\\
v_k^{\mathrm{full}}(b;\theta_k)&=\min_{c(a)\le b}\;p_k(a;\theta_k),
\end{align}
with $p_k$ the starvation probability of focal category $k$ and the minima
over prior-available actions. Let $\mathcal{B}_k$ be the finite set of action
costs of category $k$ not exceeding the plan budget $B$. Writing $\ubound$
for the bound, to keep it distinct from that budget,
\begin{equation}
\evpi\;\le\;\ubound\;:=\;\frac{1}{K_0}\sum_k I_k,
\qquad
I_k=\mathbb{E}_{\theta_k}\!\left[\max_{b\in\mathcal{B}_k}
\left(v_k^{\mathrm{inv}}(b)-v_k^{\mathrm{full}}(b;\theta_k)\right)^{\!+}
\right].
\label{eq:premium}
\end{equation}
\end{theorem}

\begin{proof}
(S1) Every plan induces a budget split by setting $\beta_k=c(S_k)$, and every
split $\{\beta_k\}$ with $\sum_k\beta_k\le B$, together with per-category
actions of cost at most $\beta_k$, induces a feasible plan. The
correspondence is not one-to-one, since several splits induce the same plan,
but that is not needed below.

\noindent(S2) By independence of the $\theta_k$ and additivity of the loss
\eqref{eq:loss}, the prior objective separates across categories once the
split is fixed.

\noindent(S3) Hence, with $A=\min_S\mathbb{E}_\theta j(S;\theta)$,
$K_0A=\min_{\text{split}}\sum_k v_k^{\mathrm{inv}}(\beta_k)$.

\noindent(S4) Similarly, with $M(\theta)=\min_S j(S;\theta)$,
$K_0M(\theta)=\min_{\text{split}}\sum_k
v_k^{\mathrm{full}}(\beta_k;\theta_k)$.

\noindent(S5) Define the pointwise deficit
$d_k(b;\theta_k)=v_k^{\mathrm{inv}}(b)-v_k^{\mathrm{full}}(b;\theta_k)$; it
may be negative for particular $\theta_k$, and only its positive part enters.

\noindent(S6) Both value functions are nonincreasing step functions of $b$
whose jumps occur only at action costs, so $d_k(\cdot;\theta_k)$ is
piecewise constant and $\sup_{b\le B}d_k(b;\theta_k)^+
=\max_{b\in\mathcal{B}_k}d_k(b;\theta_k)^+=:D_k(\theta_k)$. This is what
makes the premium finite to evaluate: a supremum over a continuum of budgets
is attained on a finite breakpoint set.

\noindent(S7) Pointwise in $\theta$ and for every split,
$\sum_k v_k^{\mathrm{full}}(\beta_k;\theta_k)\ge
\sum_k v_k^{\mathrm{inv}}(\beta_k)-\sum_k D_k(\theta_k)$.

\noindent(S8) Minimising both sides over splits preserves the inequality,
since $f\ge g-c$ pointwise implies $\min f\ge\min g-c$ for a constant $c$;
the correction term does not depend on the split, which is what allows the
minimisation to pass through. Hence
$K_0M(\theta)\ge K_0A-\sum_k D_k(\theta_k)$.

\noindent(S9) Taking $\mathbb{E}_\theta$ of (S8) and using the definition of
$I_k$ in \eqref{eq:premium} gives $K_0\,\mathbb{E}[M]\ge K_0A-\sum_k I_k$.
The slack in the bound is introduced earlier, at (S6) and (S7), where a
category's deficit is replaced by its maximum over budgets and the positive
part discards the cases in which identification would have hurt.

\noindent(S10) Rearranging, $\evpi=A-\mathbb{E}[M]\le K_0^{-1}\sum_k I_k$.
\end{proof}

\begin{corollary}[Dilution]\label{cor:dilution}
Call focal category $k$ \emph{isolated} if its subtree contains exactly one
latent category $(v,k')$, the candidates in $C_{v,k'}$ have no inputs and
supply no category other than $(v,k')$, and every other node of the subtree
fails with probability zero. Let $c_v$ be the least cost of securing a node
through which $v$'s output reaches the focal node, $v$ included, with
$c_v=\infty$ if there is none, as when $v=0$. If focal category $k$ is
isolated, with $m$ equal-cost, equal-risk candidates of risk $\bar q$, a
uniform prior and securing cost $c_u$, then
$I_k=\bar q(1-1/m)$ if $c_u\le B$ and $c_u<c_v$, and $I_k=0$ otherwise.
\end{corollary}

\begin{proof}
(D1) Isolation fixes the starvation probability. With supplier $\theta$,
$Z_k=F_\theta$ unless the plan secures $\theta$ or a node through which $v$'s
output reaches the focal node, in which case $Z_k=0$: the other candidates
feed nothing, and every other node of the subtree is reliable. Hence
$p_k(a;\theta)=0$ if $a$ contains $\theta$ or such a node, and
$p_k(a;\theta)=\bar q$ otherwise. Securing any other node changes nothing.

\noindent(D2) If $b<\min\{c_u,c_v\}$ no useful action is affordable, so both
value functions equal $\bar q$ and the deficit is zero. If
$b\ge\min\{c_v,mc_u\}$ both planners can reach zero, by the path cut or by
securing every candidate, and the deficit is zero.

\noindent(D3) Suppose $c_u<c_v$ and
$b\in[jc_u,(j+1)c_u)\cap[c_u,\min\{c_v,mc_u\})$ with $1\le j<m$. The
informed planner secures the true candidate, so $v^{\mathrm{full}}=0$. The
uninformed planner can afford at most $j$ candidates; under the uniform prior
any $j$ of them cover the supplier with probability $j/m$, so
$v^{\mathrm{inv}}=\bar q(m-j)/m$. The deficit $\bar q(m-j)/m$ is largest at
$j=1$, on $[c_u,\min\{2c_u,c_v\})$, which is nonempty because $c_u<c_v$ and
$c_u>0$.

\noindent(D4) That interval contains the breakpoint $c_u$, which lies in
$\mathcal B_k$ exactly when $c_u\le B$. Then $D_k(\theta)=\bar q(1-1/m)$ for
every $\theta$ and $I_k=\bar q(1-1/m)$. If $c_u>B$, every breakpoint in
$\mathcal B_k$ falls under (D2). If $c_u\ge c_v$, every budget is either below
$c_v\le c_u$ or at least $c_v$, and again falls under (D2). In both cases
$I_k=0$.
\end{proof}

Isolation is not a formality. In Example~\ref{ex:marginal} of
Appendix~\ref{app:H1}, raise the cost of securing $w$ above $c_u$ and keep
$c_u\le B$. Focal category $1$ then has two equal-risk candidates with
$\bar q=0.1$, but $u_1$ also supplies $(w,b)$, so the category is not
isolated. Its premium is $I_1=0.025$, not $\bar q(1-1/m)=0.05$: when $u_2$ is
the supplier, securing it still leaves $w$ exposed through $u_1$, and the
informed planner's advantage vanishes.

\paragraph{When the bound is attained.} Corollary~\ref{cor:dilution} gives
the premium $I_k$ exactly. It does not say that $\evpi=\ubound$, and in
general the two differ. Remark~\ref{rem:loose} gives one reason, an
affordable path cut. A second is that the uninformed planner can move budget
between categories. Take two isolated categories, each with $m=2$ candidates
and $c_u=1$, risks $0.1$ and $0.9$, unaffordable path cuts and $B=2$. The
uninformed planner secures both candidates of the riskier category, so
$\evpi=0.05$ while $\ubound=0.25$, although $c_u\le B<c_v$ in both
categories. Equality does hold in the following case. Suppose the
assumptions of Theorem~\ref{thm:premium} hold, the $K_0$ focal categories
are isolated with a common $m$, $\bar q$ and $c_u$ and a uniform prior, every
path cut costs more than $B$, and $K_0c_u\le B<(K_0+1)c_u$. Then
$\evpi=\ubound=\bar q(1-1/m)$.

\noindent(E1) Every category has $c_u\le B<c_v$, so
Corollary~\ref{cor:dilution} gives $I_k=\bar q(1-1/m)$ for each $k$, and
$\ubound=\bar q(1-1/m)$.

\noindent(E2) Securing the $K_0$ true suppliers costs $K_0c_u\le B$ and sets
every $Z_k$ to zero, so $\min_Sj(S;\theta)=0$ for every $\theta$.

\noindent(E3) A plan the uninformed planner can afford contains no path cut
and at most $\lfloor B/c_u\rfloor=K_0$ candidates, and securing any other
node of the subtrees has no effect, by (D1). If it secures $n_k$ candidates
of category $k$, that category starves with prior probability
$\bar q(m-n_k)^+/m$. Since $\sum_kn_k\le K_0$, $K_0$ times its objective is
at least $\bar q\sum_k(m-n_k)/m\ge\bar qK_0(m-1)/m$, with equality at
$n_k=1$ for every $k$.

\noindent(E4) Hence the prior-optimal value is $\bar q(1-1/m)$, and by (E2)
$\evpi=\bar q(1-1/m)=\ubound$.

\begin{proposition}[A checkable witness for (H2)]\label{prop:witness}
Fix a latent category $(v,k)$ with $v\ne0$, so that $v$ is itself securable,
and suppose that (i) every candidate $u\in C_{v,k}$ has $v$ as its only
consumer, and (ii) $c_v\le\min_{u\in C_{v,k}}c_u$. Then securing $v$ is an identification-free replacement
satisfying both clauses of (H2) for every $u\in C_{v,k}$. Both conditions are
functions of the public graph and the public cost vector, so they can be
checked before any structural discovery is commissioned.
\end{proposition}

\begin{proof}
(W1) Securing $v$ guarantees $v$'s output to its consumers in every scenario
and for every $\theta_{v,k}$, so under any plan containing that action the
focal loss is a function of the scenario and of the plan's other elements
alone, and its conditional law does not depend on $\theta_{v,k}$: the action is
identification-free in the sense of Definition~\ref{def:free}. The service of
$(v,k)$ itself is unchanged; what the action removes is its influence on $L$.

\noindent(W2) The cost clause is (ii) directly: $c_v\le c_u$ for every
candidate $u\in C_{v,k}$.

\noindent(W3) For the dominance clause, fix $\theta$ and a plan $S$ containing
$u$. Under the conjunctive cascade, securing $u$ affects the focal loss only
through the categories $u$ serves, and by (i) the only such category is
$(v,k)$; its effect there is to remove $u$'s own failure from that category,
which helps only when $u$ is the true supplier and would otherwise fail.
Securing $v$ instead guarantees $v$'s output outright, which removes the
failure of $(v,k)$ and of every other category $v$ consumes, for every
$\theta$ and every scenario.

\noindent(W4) Hence the operability of every node downstream of $v$ under
$S\setminus\{u\}\cup\{v\}$ is at least that under $S$, scenario by scenario,
so $L(S\setminus\{u\}\cup\{v\};\theta,\omega)\le L(S;\theta,\omega)$
pointwise, and therefore $j(S\setminus\{u\}\cup\{v\};\theta)\le j(S;\theta)$
for every $\theta$. With (W1) and (W2), (H2) holds with $a(u)=v$.
\end{proof}

The hypothesis $v\ne0$ is not cosmetic. Plans are subsets of
$V\setminus\{0\}$, so the focal node cannot be secured and the witness is
unavailable for exactly the $K_0$ categories that enter the loss directly.
Proposition~\ref{prop:witness} closes the gap left by (H2) in one direction
only. It gives a sufficient condition a firm can test on its own graph and
cost data; it does not decide (H2) in general, and its condition (i) fails
wherever a candidate supplies more than one consumer.
Section~\ref{sec:limits} records what remains open.

\begin{remark}[The premium is loose where the theorem bites]\label{rem:loose}
The maximum in \eqref{eq:premium} ranges over all breakpoints, including
budgets a planner would not choose. If $c_v\le B$, so that the consumer cut is
affordable, both planners take it and reach zero loss for that category, so
its contribution to EVPI is zero; this follows directly, not through
Theorem~\ref{thm:zero}, whose cost clause does not hold when $c_u<c_v$. Yet
Corollary~\ref{cor:dilution} still returns $I_k=\bar q(1-1/m)>0$ in exactly
that case. The bound remains valid, being an upper bound, but it is not tight
there, and a positive premium therefore does not show that identifying the
channel has any decision value. This accounts for a substantial part of the mass near zero in
Figure~\ref{fig:tight}: among the enumerated generic instances in which every
category has an affordable cut, the mean tightness ratio is $0.102$ and
$67\%$ of instances fall below $0.05$, against $0.203$ and $41\%$ among those
in which no category does. The looseness is a property of the bound itself.
\end{remark}

\begin{corollary}[Per-channel bound]\label{cor:evppi}
Under the assumptions of Theorem~\ref{thm:premium}, the partial value of
learning $\theta_k$ alone satisfies $\evppi_k\le I_k/K_0$, where
$\evppi_k=A-\mathbb{E}_{\theta_k}\!\left[\min_S\mathbb{E}_{\theta_{-k}}
j(S;\theta)\right]$.
\end{corollary}

\begin{proof}
(P1) With $\theta_k$ known and the remaining components unknown, the
separation of step (S2) gives, for each split,
$K_0\,\mathbb{E}_{\theta_{-k}}j=v_k^{\mathrm{full}}(\beta_k;\theta_k)
+\sum_{j\ne k}v_j^{\mathrm{inv}}(\beta_j)$.

\noindent(P2) By (S6), $v_k^{\mathrm{full}}(\beta_k;\theta_k)\ge
v_k^{\mathrm{inv}}(\beta_k)-D_k(\theta_k)$ for every $\beta_k$.

\noindent(P3) Adding $\sum_{j\ne k}v_j^{\mathrm{inv}}(\beta_j)$ to both sides
of (P2) and minimising over splits, permissible because $D_k(\theta_k)$ does
not depend on the split, gives $K_0\,\min_S\mathbb{E}_{\theta_{-k}}j\ge
K_0A-D_k(\theta_k)$.

\noindent(P4) Taking $\mathbb{E}_{\theta_k}$ and rearranging gives
$\evppi_k\le I_k/K_0$.
\end{proof}

\subsection{Dispersion under a specified construction}\label{sec:disp}

Heterogeneous latent candidate risk can be introduced through a dispersion
parameter $\lambda$. Monotonicity in a vague notion of ``more dispersion''
would not be well posed, since it leaves open which ordering of distributions
is meant, so we fix the family and compute the premium exactly.

\begin{definition}[Two-point dispersion family $\mathcal D(\lambda)$]
\label{def:disp}
A latent category has $m$ candidates and a uniform prior over which of them
is the supplier. Candidate risks take the two values $\bar q\pm\lambda\delta$
with equal probability, independently of which candidate is the supplier,
where $\lambda\in[0,1]$, $\delta>0$ and $\bar q\pm\delta\in[0,1]$. The
assignment of risks to candidates is latent and exchangeable under the
prior, and is revealed only by probing. The category is an isolated focal
category in the sense of Corollary~\ref{cor:dilution}, and $c_v$ is the
least cost of a path cut as defined there. Actions are: no action at cost $0$;
secure a named candidate at cost $c_u$; secure the consumer at cost
$c_v>c_u$.
\end{definition}

\begin{proposition}[Premium under $\mathcal D(\lambda)$]\label{prop:disp}
Write $D=\bar q(1-1/m)$ for the dilution term of
Corollary~\ref{cor:dilution}. Under $\mathcal D(\lambda)$,
\begin{equation}
I_k(\lambda)=
\begin{cases}
\tfrac12\max\{\lambda\delta,\,D\}+\tfrac12 D, & c_u\le B,\\[3pt]
\tfrac12\lambda\delta, & c_u>B .
\end{cases}
\label{eq:dispprem}
\end{equation}
In both cases $I_k$ is nondecreasing in $\lambda$. In the first it is
constant for $\lambda\le D/\delta$ and strictly increasing above that
threshold; in the second it is strictly increasing throughout. If the risk
assignment is instead public, the prior planner can target the fragile
candidate directly and the premium need not increase in $\lambda$.
\end{proposition}

\begin{proof}
(R1) At $b=0$ the only available action is inaction, so
$v^{\mathrm{inv}}(0)=\mathbb{E}[q_\theta]=\bar q$ by the mean-preserving
construction, while $v^{\mathrm{full}}(0;\theta)=q_\theta$. The deficit is
$\bar q-q_\theta$, whose positive part equals $\lambda\delta$ when the
supplier carries the low risk and $0$ otherwise, each with probability
$\tfrac12$.

\noindent(R2) For $b\in[c_u,c_v)$ the informed planner secures the supplier,
so $v^{\mathrm{full}}=0$. The uninformed planner faces exchangeable
candidates --- this is where latency of the assignment is used --- so every
named candidate has the same prior value; securing one leaves the category
exposed with probability $(m-1)/m$ at conditional risk $\bar q$, giving
$v^{\mathrm{inv}}=D$. This is below the do-nothing value $\bar q$, so it is
the prior planner's best action on that interval. The deficit is $D$ for
both risk types.

\noindent(R3) For $b\ge c_v$ both planners cut the consumer and reach zero,
so the deficit is zero.

\noindent(R4) If $c_u\le B$, the maximum over breakpoints is
$\max\{\lambda\delta,D\}$ when the supplier carries the low risk, and $D$
when it carries the high risk. Averaging over the two equally likely types
gives the first line of \eqref{eq:dispprem}.

\noindent(R5) If $c_u>B$, only the breakpoint $b=0$ is feasible and (R1)
gives the second line.

\noindent(R6) Both expressions are nondecreasing in $\lambda$ because
$\lambda\mapsto\max\{\lambda\delta,D\}$ and $\lambda\mapsto\lambda\delta$
are nondecreasing; the threshold statement follows from
$\lambda\delta\ge D\iff\lambda\ge D/\delta$.

\noindent(R7) If the assignment is public, the prior planner may secure the
high-risk candidate deliberately, which changes $v^{\mathrm{inv}}$ and
removes the exchangeability used in (R2), so monotonicity no longer follows
from this argument.
\end{proof}

Equation \eqref{eq:dispprem} predicts a threshold rather than a slope. Below
$\lambda=D/\delta$ the dilution term dominates and dispersion contributes
nothing to the premium; only above it does dispersion matter.

\subsection{Why the bound does not transfer to CVaR}\label{sec:cvarhonest}

Theorem~\ref{thm:zero} extends to CVaR at no cost
(Corollary~\ref{cor:lawinv}). Theorem~\ref{thm:premium} does not, and the
obstruction is structural.

\begin{proposition}[What can be said]\label{prop:cvartransfer}
Let $L\in[0,1]$. Then $\mathbb{E}[L]\le\cvar_\alpha(L)\le
\mathbb{E}[L]/(1-\alpha)$, and, writing $\evpi_{\cvar}$ for the value of
perfect information under the CVaR objective,
\begin{equation}
\evpi_{\cvar}\;\le\;\frac{1}{1-\alpha}\,\evpi
\;+\;\frac{\alpha}{1-\alpha}\,\mathbb{E}_\theta\!\left[\min_S
j(S;\theta)\right].
\label{eq:cvarbound}
\end{equation}
\end{proposition}

\begin{proof}
(V1) By the Rockafellar--Uryasev representation~\cite{rockafellar2000},
$\cvar_\alpha(L)=\min_t\{t+(1-\alpha)^{-1}\mathbb{E}(L-t)^+\}$.

\noindent(V2) Evaluating at $t=0$ and using $L\ge0$ gives
$\cvar_\alpha(L)\le\mathbb{E}[L]/(1-\alpha)$.

\noindent(V3) For any $t$, since $(1-\alpha)^{-1}\ge1$ and
$\mathbb{E}(L-t)^+\ge0$,
$t+(1-\alpha)^{-1}\mathbb{E}(L-t)^+\ge t+\mathbb{E}(L-t)^+
\ge t+\mathbb{E}(L-t)=\mathbb{E}[L]$; minimising over $t$ gives
$\cvar_\alpha(L)\ge\mathbb{E}[L]$.

\noindent(V4) With $A_{\cvar}=\min_S\mathbb{E}_\theta J(S;\theta)$, (V2)
gives $A_{\cvar}\le(1-\alpha)^{-1}A$.

\noindent(V5) By (V3),
$\mathbb{E}_\theta\min_S J(S;\theta)\ge\mathbb{E}_\theta\min_S j(S;\theta)
=\mathbb{E}[M]$.

\noindent(V6) Subtracting and using $A=\evpi+\mathbb{E}[M]$ yields
\eqref{eq:cvarbound}.
\end{proof}

Inequality \eqref{eq:cvarbound} is not a multiplicative transfer. Its second
term is proportional to the residual loss $\mathbb{E}_\theta[\min_S
j(S;\theta)]$, the expected loss that remains after optimal mitigation, which
is of the same order as the objective itself. At $\alpha=0.9$ the coefficient
$\alpha/(1-\alpha)$ is $9$, so that term is nine times a quantity of the order
of the objective, and the inequality is too weak to be useful. We therefore
do not use it. The obstruction is visible at step (S2) of
Theorem~\ref{thm:premium}, which requires the objective to be \emph{additive}
across categories. CVaR is subadditive, and subadditivity gives an
inequality, the second statement of Proposition~\ref{prop:bern}, in the
direction that destroys the decomposition. A CVaR version
of Theorem~\ref{thm:premium} would need a different argument, which we do
not have.

\subsection{Unresolved extension: shared subtrees}\label{sec:shared}

When category subtrees share nodes, step (S2) fails and the decomposition
fails with it. We state what we believe and have not proved.

\begin{conjecture}\label{conj:shared}
Under a single-cut regime, meaning that optimal plans take at most one
action per focal-category channel, and with replacement slack, the bound
\eqref{eq:premium} holds when each shared channel's premium is counted once
per focal category it feeds.
\end{conjecture}

\noindent Section~\ref{sec:verify} reports an enumeration search for
counterexamples that found none, which is evidence about a finite class of
small instances and not a proof.

\begin{table}[ht]
\caption{Assumptions, results and where each applies.}\label{tab:applic}
\small
\begin{tabular}{@{}p{2.2cm}p{3.8cm}p{3.1cm}p{2.7cm}@{}}
\toprule
result & required assumptions & establishes & applicability \\
\midrule
Lem.~\ref{lem:garbling} & finite feasible plan set &
acquisition value $\le\evpi$ & general \\
Prop.~\ref{prop:H} & implemented planner defined in both regimes &
$H=\evpi+\varepsilon_\pi-\varepsilon_\theta$ & general \\
Thm.~\ref{thm:zero} & (H1), (H2) & $\evpi=0$ & (H1) via (H1$'$), checkable in advance;
(H2) via Prop.~\ref{prop:witness} or a plan-space argument \\
Prop.~\ref{prop:witness} & $v\ne0$, candidates sole-consumed by $v$,
$c_v\le\min_u c_u$ & (H2) holds for that category & checkable from public
data \\
Cor.~\ref{cor:evppi} & as Thm.~\ref{thm:premium} &
$\evppi_k\le I_k/K_0$ & per category \\
Cor.~\ref{cor:lawinv} & (H1), (H2) for a monotone law-invariant objective &
zero value under CVaR & theoretical \\
Thm.~\ref{thm:premium} & disjoint subtrees, independent $\theta_k$,
additive budget & $\evpi\le\ubound$ & subset of networks \\
Cor.~\ref{cor:dilution} & isolated category, equal cost and risk, uniform prior &
exact premium $I_k$; $\evpi=\ubound$ only in a sub-case & special case \\
Prop.~\ref{prop:disp} & family $\mathcal D(\lambda)$ &
premium in closed form & specified family \\
Conj.~\ref{conj:shared} & shared channels, single cut, slack &
unproved bound & exploratory \\
\bottomrule
\end{tabular}
\end{table}

Table~\ref{tab:applic} collects the assumptions each result needs and where
it applies, since the hypotheses differ from result to result and the
distinctions matter for what may be claimed in any application.

\section{Computational verification}\label{sec:verify}

The experiments in this section are not offered as proof of the analytical
results. They serve two purposes: verifying that the implementation matches
the definitions, and searching finite instances for violations that would
reveal errors in the stated results. In each category a plan may take no
action, cut at the consumer, or secure any nonempty subset of the
candidates. Every instance is enumerated exactly, over all plans and all $\theta$, with no
sampling inside an instance; seeds are used only to draw the instances
themselves, and the seeds are listed in the statement on data and code
availability.

\begin{table}[ht]
\caption{Enumeration checks. A single violation would falsify the
corresponding statement.}\label{tab:verify}
\small
\begin{tabular}{@{}lrrlr@{}}
\toprule
batch & inst. & viol. & auxiliary check & tightness \\
\midrule
generic random & 6000 & 0 & --- & mean 0.132, max 0.985 \\
Thm.~\ref{thm:zero} regime & 2000 & 0 & nonzero-EVPI cases: 0 & --- \\
Cor.~\ref{cor:dilution} regime & 2000 & 0 & formula mismatches: 0 & mean 0.212, max 1.000 \\
\bottomrule
\end{tabular}
\end{table}

Table~\ref{tab:verify} reports the three batches, and Figure~\ref{fig:tight}
shows the distribution of the ratio of exact EVPI to the bound. Most instances lie near zero, so the bound is loose in general.
The dilution batch is loose too, with mean ratio $0.212$: the premium is
exact there, but the bound is not. Equality is reached only on some of its
instances, and it holds exactly in the sub-case identified after
Corollary~\ref{cor:dilution}, so the bound cannot be multiplied by any
constant smaller than one without further assumptions.

\begin{figure}[ht]
\centering
\includegraphics[width=0.78\textwidth]{fig4_tightness.png}
\caption{Ratio of exact EVPI to the bound of Theorem~\ref{thm:premium} for
all 6000 enumerated generic instances, every one of which is plotted. Ratios
within $10^{-12}$ of zero, floating-point residues of an exact zero, are
plotted as zero. No value exceeds one.}
\label{fig:tight}
\end{figure}

\paragraph{The premium against exact EVPPI.} On miniature disjoint instances
we computed the exact per-channel $\evppi_k$ by enumeration and compared it
with the quantity Corollary~\ref{cor:evppi} says bounds it, namely
$I_k/K_0$. Both sides carry the $1/K_0$ factor, so the slack reported here is
slack in the bound and not the divisor. Over five independent seeds and
$24{,}981$ channels, $I_k/K_0$ never fell below the enumerated $\evppi_k$,
and the whole-instance bound was never violated. On $25.6\%$ of channels
$\evppi_k$ is zero while the premium $I_k/K_0$ is positive, so on those
channels the bound carries no information about the value of identification.
Where $\evppi_k$ is strictly positive, the premium is a median $1.94$ times
it, with interquartile range $1.37$ to $3.50$. Over all channels the median
absolute gap $I_k/K_0-\evppi_k$ is $0.023$, with interquartile range $0.011$
to $0.040$. Rank agreement is moderate and stable across seeds (Spearman
$0.501$ to $0.518$; Pearson $0.573$ to $0.589$). The rank statistic is
computed after merging values within $10^{-12}$ of each other, because on
about a quarter of channels, the $25.6\%$ above, $\evppi_k$ is
mathematically zero but stored with floating-point residue, which would
otherwise make its fourth decimal arbitrary.

\paragraph{Search for a counterexample to Conjecture~\ref{conj:shared}.}
Over five seeds and $100{,}000$ instances in which channels feed several
focal categories, the conjectured bound was never violated and the minimum
slack was zero, so it is attained with equality on some instances. The
search covers instances with at most three channels and three focal
categories; a counterexample may exist outside that range.

\section{Scope and limitations}\label{sec:limits}

(1) The conditions and the bound are not evaluated on real supply networks
here. Every result is either proved or checked by exact enumeration on
synthetic miniature instances. Whether (H1$'$) or (H2) holds on any observed
network, how large the identification premiums are there, and how close the
bound comes to $\evpi$ at realistic scale are open questions. Answering them
on observed multi-tier structure, with the latent layers it actually has, is
the natural next step and is left to future work.

(2) The verification is evidence about a finite class of small instances,
with at most three focal categories and three candidates per category, and
is not proof. A violation outside that class is not excluded by it.

(3) Theorem~\ref{thm:premium} requires node-disjoint categories. Real
networks typically share upstream nodes, and Conjecture~\ref{conj:shared},
which would cover them, is open; the enumeration search supports it only on
instances with at most three channels and three focal categories. The CVaR
version of Theorem~\ref{thm:premium} does not follow from the expected-loss
version, and Proposition~\ref{prop:cvartransfer} is too weak to substitute
for it.

(4) The bound can be loose. Remark~\ref{rem:loose} identifies one source, an
affordable consumer cut, and on the generic enumerated instances the mean
ratio of $\evpi$ to the bound is $0.132$. It is attained with equality only
in the sub-case identified after Corollary~\ref{cor:dilution}.

(5) (H2) has no general verification procedure. Proposition~\ref{prop:witness}
is one sufficient route; its condition (i) fails wherever a candidate
supplies more than one consumer, and in the family of
Definition~\ref{def:disp} the cost ordering makes (H2) fail.

(6) A firm that runs a heuristic planner with and without $\theta$ observes
$H$, not $\evpi$. Proposition~\ref{prop:H} relates the two by an identity
whose sign is undetermined, and we give no bound on the gap between them.

(7) All results concern one mitigation problem: budgeted backup
qualification under a conjunctive cascade. Inventory buffering, dual
sourcing and contractual remedies may value structural information
differently. Concurrent dual sourcing in particular is outside the model, as
Section~\ref{sec:model} notes: the candidate set is epistemic, and no
category is served by two suppliers at once.

\section{Conclusion}\label{sec:concl}

When the candidates of every supply category are exchangeable by a symmetry
of the public network that carries their consumers, upstream subtrees and
correlation clusters, \emph{and} identification-free actions of no greater cost uniformly dominate, the
decision value of hidden supply-network structure is exactly zero, and by Lemma~\ref{lem:garbling} no
information-acquisition policy can recover it. The statement concerns the
decision problem rather than any method, and it holds for any monotone
law-invariant objective, including CVaR. Where those conditions fail, the
value is bounded above, for node-disjoint categories, by identification
premiums computable in advance from prior and public data, and the bound is attained in a sub-case of a dilution special case.
Exact enumeration of synthetic miniature instances found no violation of any
of these statements and no counterexample to the conjectured extension to
shared subtrees; it also shows the bound to be loose on typical instances.
None of this has been evaluated on a real supply network.

The practical reading is that the value can be examined in advance. Before commissioning sub-tier discovery, a
firm can ask whether its situation satisfies conditions under which the
answer is provably nothing, and if not, what the largest possible payoff
would be. Verifying (H1) through (H1$'$) is a search for automorphisms of the
public structure, not a comparison of the candidates' marginal risks, costs or
prior weights, and it is global: a single category whose candidates differ in
their consumers, upstream subtrees or clusters defeats it.
Verifying (H2) requires an argument over the plan space, for which we give no
general procedure, and its cost clause is a real restriction rather than a
formality. Within the model there is one route to it: a consumer-level
securing action that is no dearer than securing a named candidate.
Proposition~\ref{prop:witness} makes that route checkable from public data
when, in addition, the consumer is not the focal node and every candidate has
it as its only consumer.
The cost condition is exactly what Definition~\ref{def:disp} makes false:
there the consumer cut costs $c_v>c_u$ and securing every candidate costs
$mc_u$, so both canonical identification-free replacements are dearer than
the identification-dependent action they replace, and (H2) fails. We have no
second in-model witness. Dual-sourcing or buffering a shared downstream stage
would serve, but the model has no inventory and no concurrently
multi-sourced category (Section~\ref{sec:model}); a contractual remedy priced
per category rather than per supplier would also serve, but it is not an
element of $V$. Either would require extending the model. The enumerated
batch of Table~\ref{tab:verify} instantiates (H1) and (H2) by construction
at miniature scale, by drawing the consumer cut below the candidate cost. No
real network is shown here to satisfy them.
Where no such action exists, the theorem is silent and
Theorem~\ref{thm:premium} is the relevant statement.

\section*{Declarations}

\paragraph{Competing interests.} The author declares no competing interests.

\paragraph{Funding.} This research received no funding from any agency in the
public, commercial or not-for-profit sectors.

\section*{Data and code availability}

No external data are used. Every reported number comes from exact
enumeration of synthetic miniature instances drawn from fixed seeds. The
code, the console output and result files of every reported run, the figure,
a specification of the computing environment and checksums of every file are
archived at Zenodo, \url{https://doi.org/10.5281/zenodo.23854894}, under the
MIT License. A single script regenerates all outputs, and the archive's
README maps each table, figure and number in the paper to the script and
output that produce it.

\begin{appendices}

\section{The sufficient condition for (H1)}\label{app:H1}

Throughout, $F_v$ is the failure indicator of node $v$, operability is
defined recursively by $O_v=1$ if $v\in S$ and
$O_v=(1-F_v)\prod_{k\in K_v}O_{s_{v,k}(\theta)}$ otherwise, where
$s_{v,k}(\theta)$ is $\theta_{v,k}$ for a latent category and the public
supplier otherwise, and $Z_k=1-O_{s_{0,k}(\theta)}$.

\begin{example}[Marginal exchangeability is not sufficient]\label{ex:marginal}
The focal node takes category $1$ from $w$ and category $2$ from $x$, both
public, so $K_0=2$. Node $w$ consumes a latent category $a$ with
$C_{w,a}=\{u_1,u_2\}$ and a public category $b$ supplied by $u_1$; $x$, $u_1$
and $u_2$ have no inputs. Failures are independent, with
$q_{u_1}=q_{u_2}=0.1$, $q_x=0.15$ and $q_w=0$. All securing costs equal $1$,
the budget is $B=1$, and the prior is uniform. The candidates have the same
failure law, cost and prior weight, and differ only in that $u_1$ also
supplies $(w,b)$.

\noindent(X1) Under the empty plan $Z_2=F_x$. If $\theta=u_1$ then
$O_w=1-F_{u_1}$, so $\mathbb P(Z_1=1)=0.1$ and $L$ takes the values
$0,\tfrac12,1$ with probabilities $0.765,0.220,0.015$. If $\theta=u_2$ then
$O_w=(1-F_{u_1})(1-F_{u_2})$, so $\mathbb P(Z_1=1)=0.19$ and the
probabilities are $0.6885,0.2830,0.0285$. The law of $L$, and even its mean
($0.125$ against $0.170$), depends on $\theta$, so (H1) fails.

\noindent(X2) Proposition~\ref{prop:witness} applies to $(w,a)$: $w\ne0$,
$w$ is the only consumer of $u_1$ and of $u_2$, and $c_w\le\min_uc_u$.
Securing $w$ sets $Z_1=0$ in every scenario, so it dominates securing either
candidate pointwise at no greater cost, and (H2) holds.

\noindent(X3) The feasible plans are $\emptyset,\{w\},\{x\},\{u_1\},\{u_2\}$,
with $j(\cdot;u_1)=0.125,0.075,0.050,0.075,0.125$ and
$j(\cdot;u_2)=0.170,0.075,0.095,0.125,0.125$. The state-optimal values are
$0.050$ and $0.075$, with mean $0.0625$; the prior-optimal plan is $\{x\}$,
with value $0.0725$. Hence $\evpi=0.01>0$. Step (T5) of
Theorem~\ref{thm:zero} fails: both state-optimal plans are
identification-free, but $j(\{x\};\theta)$ varies with $\theta$.

Variants in which $u_1$ alone has an upstream supplier that also feeds $x$,
or in which $u_1$ alone shares a copula cluster with $x$, keep the mean of
$L$ fixed and change its law, so they defeat Corollary~\ref{cor:lawinv}
while leaving expected loss unchanged.
\end{example}

\begin{lemma}[Sufficiency of (H1$'$)]\label{lem:H1suff}
Let $\sigma$ be a structure automorphism. For every plan $S$ and state
$\theta$, $Z(\sigma S;\sigma\theta,\cdot)$ has the same law as
$\Pi_{\sigma}Z(S;\theta,\cdot)$, where $\Pi_\sigma$ permutes the focal
coordinates by $\sigma_0$. Consequently (H1$'$) implies (H1), and the local
form stated in Section~\ref{sec:zero} implies (H1$'$).
\end{lemma}

\begin{proof}
(Q1) Fix a failure vector $F$ and define $G$ by $G_{\sigma v}=F_v$. The map
$\sigma$ carries the supplier graph of $\theta$ onto that of $\sigma\theta$:
for a public category this is (A1), and for a latent one
$s_{\sigma v,\sigma_vk}(\sigma\theta)=\sigma(\theta_{v,k})
=\sigma(s_{v,k}(\theta))$ by the definition of $\sigma\theta$. Hence $\sigma v$
has the same depth in the graph of $\sigma\theta$ as $v$ in the graph of
$\theta$.

\noindent(Q2) We show $O_{\sigma v}(\sigma S,\sigma\theta,G)
=O_v(S,\theta,F)$ by induction on that depth. Membership agrees because
$\sigma$ is a bijection; own failure agrees because $G_{\sigma v}=F_v$; and by
(Q1) and $K_{\sigma v}=\sigma_v(K_v)$ the inputs of $\sigma v$ are the images
$\sigma(s_{v,k}(\theta))$, $k\in K_v$, which lie at smaller depth, so the
induction hypothesis equates their operability with that of
$s_{v,k}(\theta)$. Substituting into the recursion gives the claim; at depth
zero the product is empty.

\noindent(Q3) By (A0) and (Q1) the supplier of focal category $\sigma_0k$
under $\sigma\theta$ is $\sigma(s_{0,k}(\theta))$, so (Q2) gives
$Z_{\sigma_0k}(\sigma S,\sigma\theta,G)=Z_k(S,\theta,F)$ for every $F$.

\noindent(Q4) Under the one-factor Gaussian copula,
$F_v=\mathbf 1\{\sqrt{r_g}\,Y_g+\sqrt{1-r_g}\,\varepsilon_v\le
\Phi^{-1}(q_v)\}$ with $g=\kappa(v)$, $r_g$ the cluster's copula parameter,
and $(Y_g)$, $(\varepsilon_v)$ independent standard normal. By (A3),
$G_v=F_{\sigma^{-1}v}$ is the same function of $Y'_g=Y_{\eta^{-1}g}$ and
$\varepsilon'_v=\varepsilon_{\sigma^{-1}v}$ that $F_v$ is of $Y$ and
$\varepsilon$. These are re-indexings of independent standard normal families
by bijections, so $G$ has the law of $F$. With (Q3) this gives the first
statement.

\noindent(Q5) Given an identification-free $S$ and $\theta,\theta'$, take
$\sigma$ from (H1$'$). Since $\sigma S=S$ and $\sigma\theta=\theta'$, (Q4)
gives $Z(S;\theta',\cdot)\overset{d}{=}\Pi_\sigma Z(S;\theta,\cdot)$. The
loss \eqref{eq:loss} is symmetric in the focal coordinates, so
$L(S;\theta',\cdot)\overset{d}{=}L(S;\theta,\cdot)$, which is (H1).

\noindent(Q6) For the local form, list the latent categories
$\ell_1,\dots,\ell_n$ and set $\theta^{(0)}=\theta$. At step $i$ apply the
automorphism $\tau_i$ that maps $\theta^{(i-1)}_{\ell_i}$ to
$\theta'_{\ell_i}$ (the identity if they agree). It fixes $\ell_i$, so the
$\ell_i$ coordinate becomes $\theta'_{\ell_i}$; it fixes every other latent
category together with its candidates, so the other coordinates are
unchanged. After $n$ steps the state is $\theta'$. The composition
$\tau_n\circ\cdots\circ\tau_1$ satisfies (A0)--(A3), maps $\theta$ to
$\theta'$ and fixes $S$, which is (H1$'$).
\end{proof}

The requirement $\sigma S=S$ in (H1$'$) may be weakened to
$L(\sigma S;\theta',\cdot)\overset{d}{=}L(S;\theta',\cdot)$ with the same
proof, which matters when an identification-free plan meets two candidates'
upstream cones asymmetrically without changing the loss. (H1$'$) is
sufficient, not necessary: (H1) can hold by numerical coincidence without any
symmetry. Steps (T5) of Theorem~\ref{thm:zero} and (C2) of
Corollary~\ref{cor:lawinv} use only (H1), so both results hold unchanged with
(H1$'$) in place of (H1).

\end{appendices}

\begin{thebibliography}{99}

\bibitem{artstein1993}
Artstein, Z., Wets, R.J.-B.: Sensors and information in optimization under
stochastic uncertainty. Math. Oper. Res. \textbf{18}(3), 523--547 (1993)

\bibitem{artstein1994}
Artstein, Z.: Probing for information in two-stage stochastic programming and
the associated consistency. In: Mandl, P., Hu\v{s}kov\'a, M. (eds.)
Asymptotic Statistics, pp. 21--33. Physica, Heidelberg (1994)

\bibitem{artstein1999}
Artstein, Z.: Gains and costs of information in stochastic programming.
Ann. Oper. Res. \textbf{85}, 129--152 (1999)

\bibitem{artzner1999}
Artzner, P., Delbaen, F., Eber, J.-M., Heath, D.: Coherent measures of risk.
Math. Finance \textbf{9}(3), 203--228 (1999)

\bibitem{avriel1970}
Avriel, M., Williams, A.C.: The value of information and stochastic
programming. Oper. Res. \textbf{18}(5), 947--954 (1970)

\bibitem{birge1982}
Birge, J.R.: The value of the stochastic solution in stochastic linear
programs with fixed recourse. Math. Program. \textbf{24}, 314--325 (1982)

\bibitem{blackwell1953}
Blackwell, D.: Equivalent comparisons of experiments. Ann. Math. Stat.
\textbf{24}(2), 265--272 (1953)

\bibitem{brintrup2018}
Brintrup, A., Wichmann, P., Woodall, P., McFarlane, D., Nicks, E., Krechel,
W.: Predicting hidden links in supply networks. Complexity \textbf{2018},
article 9104387, 1--12 (2018)

\bibitem{caridi2014}
Caridi, M., Moretto, A., Perego, A., Tumino, A.: The benefits of supply
chain visibility: a value assessment model. Int. J. Prod. Econ.
\textbf{151}, 1--19 (2014)

\bibitem{cermeno2026}
Cerme\~{n}o, J.E., Rodas, R.: Unlocking supply chain economic value through
multi-tier supplier network visibility. Capstone report, MIT Center for
Transportation and Logistics (2026)

\bibitem{goel2006}
Goel, V., Grossmann, I.E.: A class of stochastic programs with decision
dependent uncertainty. Math. Program. \textbf{108}, 355--394 (2006)

\bibitem{golovin2011}
Golovin, D., Krause, A.: Adaptive submodularity: theory and applications in
active learning and stochastic optimization. J. Artif. Intell. Res.
\textbf{42}, 427--486 (2011)

\bibitem{hausch1983}
Hausch, D.B., Ziemba, W.T.: Bounds on the value of information in uncertain
decision problems II. Stochastics \textbf{10}, 181--217 (1983)

\bibitem{heath2024}
Heath, A., Kunst, N., Jackson, C. (eds.): Value of Information for
Healthcare Decision-Making. Chapman and Hall/CRC, Boca Raton (2024)

\bibitem{howard1966}
Howard, R.A.: Information value theory. IEEE Trans. Syst. Sci. Cybern.
\textbf{2}(1), 22--26 (1966)

\bibitem{huang1977}
Huang, C.C., Vertinsky, I., Ziemba, W.T.: Sharp bounds on the value of
perfect information. Oper. Res. \textbf{25}(1), 128--139 (1977)

\bibitem{javdani2014}
Javdani, S., Chen, Y., Karbasi, A., Krause, A., Bagnell, D., Srinivasa, S.:
Near optimal Bayesian active learning for decision making. In: Proceedings
of AISTATS, PMLR \textbf{33}, pp. 430--438 (2014)

\bibitem{jonsbraten1998}
Jonsbr{\aa}ten, T.W., Wets, R.J.-B., Woodruff, D.L.: A class of stochastic
programs with decision dependent random elements. Ann. Oper. Res.
\textbf{82}, 83--106 (1998)

\bibitem{kosasih2022}
Kosasih, E.E., Brintrup, A.: A machine learning approach for predicting
hidden links in supply chain with graph neural networks. Int. J. Prod. Res.
\textbf{60}(17), 5380--5393 (2022)

\bibitem{ma2024}
Ma, Z., Kim, Y., Linderoth, J., Luedtke, J., Matthews, L.R.:
Probing-enhanced stochastic programming. arXiv:2407.10669 (2024)

\bibitem{madansky1960}
Madansky, A.: Inequalities for stochastic linear programming problems.
Manage. Sci. \textbf{6}(2), 197--204 (1960)

\bibitem{raiffa1961}
Raiffa, H., Schlaifer, R.: Applied Statistical Decision Theory. Harvard
Business School, Boston (1961)

\bibitem{rockafellar2000}
Rockafellar, R.T., Uryasev, S.: Optimization of conditional value-at-risk.
J. Risk \textbf{2}(3), 21--41 (2000)

\bibitem{simchi2015}
Simchi-Levi, D., Schmidt, W., Wei, Y., Zhang, P.Y., Combs, K., Ge, Y.,
Gusikhin, O., Sanders, M., Zhang, D.: Identifying risks and mitigating
disruptions in the automotive supply chain. Interfaces \textbf{45}(5),
375--390 (2015)

\bibitem{snyder2016}
Snyder, L.V., Atan, Z., Peng, P., Rong, Y., Schmitt, A.J., Sinsoysal, B.:
OR/MS models for supply chain disruptions: a review. IIE Trans.
\textbf{48}(2), 89--109 (2016)

\bibitem{strong2014}
Strong, M., Oakley, J.E., Brennan, A.: Estimating multiparameter partial
expected value of perfect information from a probabilistic sensitivity
analysis sample. Med. Decis. Making \textbf{34}(3), 311--326 (2014)

\bibitem{tomlin2006}
Tomlin, B.: On the value of mitigation and contingency strategies for
managing supply chain disruption risks. Manage. Sci. \textbf{52}(5),
639--657 (2006)

\bibitem{vayanos2025}
Vayanos, P., Georghiou, A., Yu, H.: Robust optimization with
decision-dependent information discovery. Manage. Sci. \textbf{72}(2),
1509--1528 (2025)

\bibitem{wagner2010}
Wagner, S.M., Neshat, N.: Assessing the vulnerability of supply chains using
graph theory. Int. J. Prod. Econ. \textbf{126}(1), 121--129 (2010)

\bibitem{ziemba1975}
Ziemba, W.T., Butterworth, J.E.: Bounds on the value of information in
uncertain decision problems. Stochastics \textbf{1}, 361--378 (1975)

\end{thebibliography}

\end{document}
